\section{The Unresolved Problem of Capacity Utilization}

\paragraph{A}
This paper examines how the rate of profit, capital accumulation, and productive capacity formation interact under conditions of distributional conflict between labor and capital.

\paragraph{B}
\PLACEHOLDER{Section 1 is unwritten. What exists above is a single orienting sentence; the section that the locked outline requires does not yet exist. The remainder of this box states what it has to do, so the gap is legible rather than silent.}

\paragraph{C}
\PLACEHOLDER{\textbf{Burden of the section.} Establish that capacity utilization is unresolved in the treatments that have come down to us, in three distinct ways, and show that these are three faces of one problem rather than three separate difficulties. First, utilization treated as a \emph{given macro ratio}: taken as observed when its denominator is not. Second, utilization treated as an \emph{inherited benchmark}: anchored to an engineering or survey norm imported from outside the accounting framework that is supposed to give it meaning. Third, utilization treated as an \emph{empirical residual}: whatever is left after output has been fitted to capital, which presupposes the very transformation relation at issue. The third case is where the filtering literature sits, and the objection to it here is structural rather than methodological --- a filter that fits output to capital imposes a unit transformation elasticity by construction, and so conflates utilization with the object that would be needed to identify it. That argument is available from this chapter's own accounting framework and requires no external authority.}

\paragraph{D}
\PLACEHOLDER{\textbf{Bottom line up front.} Per the writing calibration in force, the contribution must be visible by the end of the second paragraph: what the paper examines, how it identifies the transformation elasticity, and what the comparison of the United States and Chile establishes. The reader should not have to reach Section 3 to learn what is being claimed.}

\paragraph{E}
\PLACEHOLDER{\textbf{What must not appear here.} Long literature catalogues. The full center--periphery discussion, which belongs in Section 5. Estimator debates, which belong in Section 4 and Appendix C. Detailed methodological caveats. The section motivates a problem; it does not litigate the solution.}

\paragraph{F}
\PLACEHOLDER{\textbf{Closing-line burden (locked).} The section must arrive at: \emph{the utilization problem cannot be resolved at the level of the ratio alone; it must be reconstructed through the relation that governs productive-capacity formation.} This hands off directly to Section 2.}

\paragraph{G}
\PLACEHOLDER{\textbf{Also unresolved: the abstract's sample window.} The abstract states a United States window of 1947--2024, while Sections 4 through 6 estimate over 1929--2024 with MATHPLACEHOLDER0XYZ. The two cannot both be right. Note further that 1929--2024 spans 96 years against a reported MATHPLACEHOLDER1XYZ of 94, and that the estimation prose elsewhere gives the effective span as 1931--2024 after lag-1 conditioning and listwise deletion --- which is consistent with MATHPLACEHOLDER2XYZ and suggests the table titles, not the sample, are what needs correcting.}